% SOURCE_AUTHORITY: sga5_fr_workpass.tex, lines 4928--4962 (LF-normalized unit).
% SOURCE_SHA256: 791F4EFFC5E02832D5D77ED03518C8156D6F07E4C8238B03545DB93D883FBB28
\section*{II. Correspondências equivariantes}

\subsection*{5. Traços não comutativos}

Fixemos neste número um topos $T$ e um anel comutativo unitário $K$ de $T$. Os anéis considerados serão supostos associativos e unitários, mas não necessariamente comutativos. Chamaremos de $K$-álgebra todo anel $A$ (associativo e unitário) de $T$ munido de um homomorfismo (unitário) $K\to A$ cuja imagem esteja contida no centro.

\subsubsection*{5.0. Introdução}

Seja $E$ um $K$-módulo que localmente é fator direto de um módulo livre de tipo finito. Dispõe-se então de um homomorfismo de traço, de natureza local,
\begin{equation}\tag{5.0.1}
\Tr_K:\uHom_K(E,E)\longrightarrow K,
\end{equation}
definido compondo a flecha de avaliação
\begin{equation}\tag{5.0.2}
\uHom_K(E,K)\otimes_KE\longrightarrow K,\qquad f\otimes x\longmapsto f(x),
\end{equation}
com o inverso do isomorfismo canônico de produto tensorial
\begin{equation}\tag{5.0.3}
\uHom_K(E,K)\otimes_KE\longrightarrow\uHom_K(E,E),\qquad f\otimes y\longmapsto (x\longmapsto f(x)y).
\end{equation}
As flechas 5.0.2 e 5.0.3 definem-se mais geralmente para todo $K$-módulo $E$, e delas se deduzem, para todo complexo de $K$-módulos $E$, homomorfismos de complexos
\begin{align}
\uHom_K^\bullet(E,K)\otimes_KE&\longrightarrow K,\qquad f\otimes x\longmapsto f(x),\tag{5.0.4}\\
\uHom_K^\bullet(E,K)\otimes_KE&\longrightarrow \uHom_K^\bullet(E,E),\qquad f\otimes y\longmapsto (x\longmapsto f(x)y).\tag{5.0.5}
\end{align}
Por derivação, 5.0.4 e 5.0.5 dão, para $E\in\ob D^-(K)$, flechas de $D(K)$
\begin{align}
\uRHom_K(E,K)\lten_KE&\longrightarrow K,\tag{5.0.6}\\
\uRHom_K(E,K)\lten_KE&\longrightarrow\uRHom_K(E,E).\tag{5.0.7}
\end{align}
Se $E$ é perfeito (SGA 6 I), 5.0.7 é um isomorfismo, como se verifica facilmente reduzindo ao caso $E=K$; compondo o inverso de 5.0.7 com 5.0.6, obtém-se um homomorfismo de traço, de natureza global,
\begin{equation}\tag{5.0.8}
\Tr_K:\uRHom_K(E,E)\longrightarrow K,
\end{equation}
que generaliza 5.0.1 (para mais detalhes, ver SGA 6 I 7, 8).
